% META: {"id": "TSPE-CONCLGN-ALG-018", "chapitre": "TSPE-PROBABILITES", "type_objet": "algorithme", "capacites_codes": ["TSPE-CONCLGN-C1"], "programme_alignment": "OFFICIAL_ALGORITHM_EXAMPLE", "mode_creation": "ex_nihilo", "sources_inspiration": [], "status": "generated"}

\section{Algorithmique : simuler une marche aléatoire}

% ============================================================
% STRATE 1 --- Le probleme, le modele, l'algorithme
% ============================================================

\margeAppui{
Le programme cite la simulation d'une marche aleatoire parmi ses
\textbf{exemples d'algorithme} de la partie « Concentration, loi des grands
nombres ». Il ne l'impose pas : cette page est un entrainement.
}

Un promeneur part de l'origine sur une droite graduee. A chaque seconde il
avance d'un pas vers la droite ou vers la gauche, avec la meme chance,
independamment des pas precedents. A quelle distance de son point de depart
se trouve-t-il au bout de $n$ pas ?

\propriete[Esperance et variance de la position]{
Notons $X_1,\dots,X_n$ les pas, a valeurs dans $\{-1\,;\,1\}$, independants,
de meme loi avec $P(X_i=1)=P(X_i=-1)=\frac{1}{2}$, et
$S_n = X_1+\cdots+X_n$ la position apres $n$ pas. Alors
\[
  E(X_i)=0,\quad V(X_i)=E(X_i^2)-E(X_i)^2 = 1-0 = 1,
\]
donc, par linearite de l'esperance et additivite de la variance de variables
independantes,
\[
  E(S_n)=0
  \qquad\text{et}\qquad
  V(S_n)=n,\quad \sigma(S_n)=\sqrt{n}.
\]
}

\margeAppui{
$E(S_n)=0$ ne signifie pas que le promeneur reste pres de $0$ : c'est la
\emph{moyenne} des positions qui est nulle, par symetrie entre la gauche et
la droite. L'ecart-type $\sqrt{n}$ donne l'ordre de grandeur de l'eloignement
reellement observe.
}

\propriete[Ce que dit l'inegalite de Bienaymé-Tchebychev]{
En appliquant l'inegalite a $Y=S_n$, d'esperance $0$ et de variance $n$, avec
$\delta = k\sqrt{n}$ pour un reel $k>0$ :
\[
  P\big(|S_n| \geq k\sqrt{n}\,\big) \leq \frac{n}{\big(k\sqrt{n}\big)^{2}}
  = \frac{1}{k^{2}}.
\]
La majoration ne depend plus de $n$ : quelle que soit la duree de la marche,
la position depasse $k$ ecarts-types avec une probabilite au plus
$\dfrac{1}{k^{2}}$.
}

\propriete[Demarche de simulation]{%
\begin{enumerate}
  \item Fixer le nombre de pas $n$, le nombre de marches $N$ et une graine.
  \item Pour chaque marche, cumuler $n$ pas tires au hasard dans
        $\{-1\,;\,1\}$ ; conserver la position finale.
  \item Pour $k=1,2,3$, compter les marches telles que
        $|S_n| \geq k\sqrt{n}$ et diviser par $N$.
  \item Comparer chaque proportion obtenue a la majoration $1/k^{2}$.
\end{enumerate}
}

\begin{python}
from math import sqrt
from random import Random

def simuler_marches(n, N, graine):
    """Renvoie les N positions finales de marches de n pas."""
    if n <= 0 or N <= 0:
        raise ValueError("n et N doivent etre strictement positifs")
    generateur = Random(graine)
    positions = []
    for _ in range(N):
        position = 0
        for _ in range(n):
            if generateur.random() < 0.5:
                position = position + 1
            else:
                position = position - 1
        positions.append(position)
    return positions

def proportion_au_dela(positions, n, k):
    """Proportion de marches dont la position finale depasse k ecarts-types."""
    seuil = k * sqrt(n)
    depassent = [s for s in positions if abs(s) >= seuil]
    return len(depassent) / len(positions)
\end{python}

\exemple{
Pour $n=100$ pas, $N=10\,000$ marches et la graine $2024$ :

\begin{center}
\begin{tabular}{|c|c|c|}
\hline
$k$ & proportion observee de $|S_{100}| \geq k\sqrt{100}$ & majoration $1/k^{2}$ \\
\hline
$1$ & $0{,}3620$ & $1$ \\
\hline
$2$ & $0{,}0579$ & $0{,}25$ \\
\hline
$3$ & $0{,}0043$ & $0{,}1111\ldots$ \\
\hline
\end{tabular}
\end{center}

Chaque proportion respecte bien sa majoration. On observe aussi que pour
$k=1$ la majoration vaut $1$ : elle est vraie, mais elle n'apprend rien. Ce
n'est qu'a partir de $k=2$ que l'inegalite devient informative.
}

% BEGIN-VERIFY
% from math import sqrt
% from random import Random
%
% def simuler_marches(n, N, graine):
%     if n <= 0 or N <= 0:
%         raise ValueError("n et N doivent etre strictement positifs")
%     generateur = Random(graine)
%     positions = []
%     for _ in range(N):
%         position = 0
%         for _ in range(n):
%             if generateur.random() < 0.5:
%                 position = position + 1
%             else:
%                 position = position - 1
%         positions.append(position)
%     return positions
%
% def proportion_au_dela(positions, n, k):
%     seuil = k * sqrt(n)
%     depassent = [s for s in positions if abs(s) >= seuil]
%     return len(depassent) / len(positions)
%
% positions = simuler_marches(100, 10000, 2024)
% assert len(positions) == 10000
% # la position finale a la parite du nombre de pas
% assert all(s % 2 == 0 for s in positions)
% assert all(-100 <= s <= 100 for s in positions)
% observees = [proportion_au_dela(positions, 100, k) for k in (1, 2, 3)]
% assert abs(observees[0] - 0.3620) < 1e-12
% assert abs(observees[1] - 0.0579) < 1e-12
% assert abs(observees[2] - 0.0043) < 1e-12
% # chaque proportion respecte la majoration de Bienayme-Tchebychev
% for k, observee in zip((1, 2, 3), observees):
%     assert observee <= 1 / k**2
% assert abs(1 / 3**2 - 0.1111) < 5e-5
% # E(S_n) = 0 et V(S_n) = n : verification exacte sur la loi d'un pas
% from fractions import Fraction
% esperance_pas = Fraction(1, 2) * 1 + Fraction(1, 2) * (-1)
% assert esperance_pas == 0
% moment2_pas = Fraction(1, 2) * 1 + Fraction(1, 2) * 1
% assert moment2_pas - esperance_pas**2 == 1
% # la majoration en k racine de n ne depend plus de n
% for n in (25, 100, 10000):
%     assert abs(n / (2 * sqrt(n))**2 - 1 / 4) < 1e-12
% for mauvais in ((0, 10), (10, 0), (-3, 4)):
%     try:
%         simuler_marches(*mauvais, 1)
%     except ValueError:
%         pass
%     else:
%         raise AssertionError("l'argument invalide aurait du etre refuse")
% END-VERIFY

% ============================================================
% STRATE 2 --- Erreurs frequentes
% ============================================================

\erreurFrequente{
\textbf{Prendre la majoration pour une approximation.} L'inegalite de
Bienaymé-Tchebychev donne un \emph{majorant}, pas une valeur approchee.
Observer $0{,}0579$ la ou elle annonce au plus $0{,}25$ ne la contredit pas :
elle est vraie pour toute loi, et cette generalite se paie par une borne
souvent large.
}

\erreurFrequente{
\textbf{Conclure de $E(S_n)=0$ que le promeneur revient pres de l'origine.}
L'esperance est nulle par symetrie : les eloignements vers la droite et vers
la gauche se compensent en moyenne. L'eloignement typique, lui, croit comme
$\sqrt{n}$ --- c'est ce que mesure l'ecart-type.
}
